% =====================================================================
%  weekly-preamble.tex — khung cho báo cáo kỹ thuật
%  Biên dịch bằng XeLaTeX (hoặc tectonic). Không dùng pdfLaTeX.
% =====================================================================
% Lề theo chuẩn văn bản: trên/dưới 25 mm, trái 30 mm (đóng gáy), phải 20 mm
\usepackage[a4paper,top=25mm,bottom=25mm,left=30mm,right=20mm,
  headheight=14pt,footskip=12mm]{geometry}
\usepackage{fontspec}
\setmainfont{texgyrepagella}[Extension=.otf, UprightFont=*-regular, BoldFont=*-bold,
  ItalicFont=*-italic, BoldItalicFont=*-bolditalic]
\setsansfont{texgyreheros}[Extension=.otf, UprightFont=*-regular, BoldFont=*-bold,
  ItalicFont=*-italic, BoldItalicFont=*-bolditalic]
\setmonofont{lmmono10}[Extension=.otf, UprightFont=*-regular, ItalicFont=*-italic,
  BoldFont=lmmonolt10-bold, Scale=0.92]
\usepackage{polyglossia}
\setdefaultlanguage{vietnamese}

\usepackage{amsmath,amssymb}
\usepackage{xcolor}
\usepackage{booktabs,array,xltabular}
\usepackage{etoolbox,needspace}
\usepackage[section]{placeins}  % hình không trôi sang mục sau
\usepackage{enumitem}
\usepackage{fancyvrb}
\usepackage{titlesec}
\usepackage{fancyhdr}
\usepackage{tcolorbox}
\tcbuselibrary{skins,breakable}
\usepackage[hidelinks,unicode]{hyperref}

\setlength{\parindent}{0pt}
\setlength{\parskip}{4pt plus 1pt}
\setlength{\emergencystretch}{3em}
\tolerance=2000
\renewcommand{\arraystretch}{1.2}
\setlength{\tabcolsep}{5pt}
% Bảng: chữ nhỏ, sát lề trái, khoảng cách trên/dưới đều nhau
\setlength{\LTleft}{0pt}\setlength{\LTright}{0pt}
\setlength{\LTpre}{6pt}\setlength{\LTpost}{6pt}
\AtBeginEnvironment{xltabular}{\small}
% Bảng ngắn: không ngắt trang giữa bảng
\newenvironment{shorttable}{\par\nopagebreak\vspace{6pt}\nopagebreak\noindent\small\tabularx}{\endtabularx\par\vspace{6pt}}
% Vị trí hình: ưu tiên ngay tại chỗ, tránh khoảng trắng lớn
\renewcommand{\topfraction}{0.9}
\renewcommand{\bottomfraction}{0.8}
\renewcommand{\textfraction}{0.07}
\renewcommand{\floatpagefraction}{0.85}
\setcounter{topnumber}{3}\setcounter{bottomnumber}{2}\setcounter{totalnumber}{4}
\setlist{itemsep=1pt,topsep=3pt,parsep=0pt}

\definecolor{brand}{HTML}{1E3A8A}
\definecolor{muted}{HTML}{6B7280}

\titleformat{\section}{\needspace{8\baselineskip}\sffamily\Large\bfseries\color{brand}}{\thesection}{0.6em}{}
\titleformat{\subsection}{\needspace{6\baselineskip}\sffamily\large\bfseries}{\thesubsection}{0.6em}{}
\titleformat{\subsubsection}{\sffamily\normalsize\bfseries}{\thesubsubsection}{0.5em}{}
\titlespacing*{\section}{0pt}{12pt}{6pt}
\titlespacing*{\subsection}{0pt}{10pt}{4pt}

\pagestyle{fancy}
\fancyhf{}
\fancyfoot[C]{\small\sffamily\thepage}
\renewcommand{\headrulewidth}{0pt}

% \code{...}: định danh trong mã, giữ khoảng trắng, ngắt dòng tại / . _ - : , = |
\ExplSyntaxOn
\NewDocumentCommand \code { m }
  {
    \texttt
      {
        \str_map_inline:nn {#1}
          {
            ##1
            \str_case:nnF {##1}
              { {/}{\allowbreak} {.}{\allowbreak} {_}{\allowbreak}
                {-}{\allowbreak} {:}{\allowbreak} {,}{\allowbreak}
                {=}{\allowbreak} {|}{\allowbreak} }
              { }
          }
      }
  }
\ExplSyntaxOff

\newcommand{\ra}{\ensuremath{\rightarrow}}
\newcommand{\src}[1]{{\small\color{muted}Mã nguồn: \code{#1}}\par}

\newcolumntype{L}[1]{>{\raggedright\arraybackslash}p{#1}}
\newcolumntype{Y}{>{\raggedright\arraybackslash}X}

% Code block nhỏ
\DefineVerbatimEnvironment{codeblock}{Verbatim}{fontsize=\small,frame=leftline,
  framerule=1.2pt,rulecolor=\color{muted},xleftmargin=6pt,framesep=6pt}

% \figslot{tên-file}{chiều cao ô}{chú thích}
%   Có figures/<tên-file>.pdf|png|jpg thì chèn hình; chưa có thì để ô trắng.
%   Tham chiếu: \ref{fig:<tên-file>}
\usepackage{graphicx,float,caption}
\captionsetup{font=small,labelfont=bf,justification=centering,skip=6pt}
\graphicspath{{figures/}}
\newcommand{\figimg}[1]{\includegraphics[width=\linewidth,height=0.6\textheight,keepaspectratio]{#1}}
\newcommand{\figslot}[3]{%
  \begin{figure}[!htbp]\centering
    \IfFileExists{figures/#1.pdf}{\figimg{#1.pdf}}{%
    \IfFileExists{figures/#1.png}{\figimg{#1.png}}{%
    \IfFileExists{figures/#1.jpg}{\figimg{#1.jpg}}{%
      {\setlength{\fboxsep}{0pt}\color{muted!60}\fbox{\parbox[b][#2][b]{\dimexpr\linewidth-0.8pt}{%
        \hfill{\scriptsize\ttfamily\color{muted!80}figures/#1\ }\par\vspace{2pt}}}}}}}
    \caption{#3}\label{fig:#1}
  \end{figure}}
